% ============================================================
\subsection{Accuracy Comparison}

\begin{table}[H]
\centering
\caption{Relative Error vs Python (double-precision reference)}
\begin{tabular}{lccc}
\toprule
Variant & Case1 & Case2 & Case3 \\
\midrule
Python (ref.) & 0 & 0 & 0 \\
V3 & $\approx 0$ & $2.32\times 10^{-7}$ & $2.35\times 10^{-5}$ \\
V4 & $\approx 0$ & $2.32\times 10^{-7}$ & $2.35\times 10^{-5}$ \\
V6 & \multicolumn{3}{c}{\textit{[to be filled]}} \\
\bottomrule
\end{tabular}
\end{table}

Numerical deviations generally increase with input size due to finite-precision effects and the accumulation of rounding errors in long summations. 
Importantly, enabling hardware multiplier support (V4) does not change the numerical outputs compared to the software baseline (V3), indicating that V4 improves performance without altering numerical correctness. 
Case 4 is omitted from this accuracy comparison because the random vector generation differs between the C implementation (\texttt{rand()}) and the Python implementation (e.g., NumPy RNG); therefore, even with the same seed, the generated input sequences are not identical and the resulting $F(x)$ values are not directly comparable.